\textbf{\textit{Soal. }}Incircle of the triangle $\triangle ABC$ is tangent to sides $BC, AC, AB$ at points $D, E, F$ respectively. Let $I$ be the incentre of triangle and $K$ be the intersection point of $AI, BC.$ Circumcircle $w$ of the triangle $\triangle IDK$ intersects $IB$ and $IC$ at points $P$ and $Q.$ Prove that lines $EF, PQ$ and the tangent to $w$ from $D$ are concurrent.

\begin{proof}[\textbf{Solusi. }] (Thursday, 20 October 2022) WLOG $AB < AC$. Misalkan $T$ adalah perpotongan kedua dari lingkaran $(FPD)$ dan $(EQD)$. Karena $BP$ dan $CQ$ masing-masing adalah garis bagi sudut $B$ dan $C$, maka $FP=PD$ dan $EQ=QD$. Lalu, karena $\dangle FPD = \dangle FTD = \dangle ETD = \dangle EQD$ maka $\triangle EQD$ dan $\triangle FPD$ sebangun. Lebih jauh, $D$ adalah pusat dari spiral similarity yang membawa $\triangle EQD$ ke $\triangle FPD$ sehingga dari well-known property, $T$ adalah perpotongan dari $EF$ dan $QP$. Oleh karena itu, akan dibuktikan bahwa $TD$ adalah garis singgung lingkaran $\omega$ di $D$.

\begin{lemma}
    $PQ \parallel BC$
    \begin{proof}[\textbf{Bukti. }]
        Perhatikan bahwa karena $\angle IKB = 180^\circ - \angle BAK - \angle KBA = 180^\circ - \frac{1}{2}\angle A - \angle B$ maka didapat $\angle PIK = 180^\circ - \angle KBI - \angle IKB =  180^\circ - \frac{1}{2}\angle B - (180^\circ - \frac{1}{2}\angle A - \angle B)= \frac{1}{2}(\angle A + \angle B) = \frac{1}{2}(180^\circ - \angle C) = 90^\circ - \frac{1}{2}\angle C$. Padahal, karena $\angle KDI=90^\circ$ didapat $IK$ adalah diameter $\omega$ sehingga $\angle KPI = 90^\circ$ yang menyebabkan $\angle IQP = \angle IKP = 90^\circ - \angle PIK = \frac{1}{2}\angle C = \angle ICB$ yang mengakibatkan $PQ \parallel BC$. Terbukti.
    \end{proof}
\end{lemma}

\begin{center}
    \begin{asy}
         /* Geogebra to Asymptote conversion, documentation at artofproblemsolving.com/Wiki go to User:Azjps/geogebra */
        import graph; size(14cm); 
        real labelscalefactor = 0.5; /* changes label-to-point distance */
        pen dps = linewidth(0.7) + fontsize(10); defaultpen(dps); /* default pen style */ 
        pen dotstyle = black; /* point style */ 
        real xmin = -7.72, xmax = 15.24, ymin = -12.22, ymax = 2.92;  /* image dimensions */
        pen zzttqq = rgb(0.6,0.2,0); pen qqwuqq = rgb(0,0.39215686274509803,0); pen ccqqqq = rgb(0.8,0,0); pen wwwwww = rgb(0.4,0.4,0.4); 
        
        draw((1.08,-0.46)--(-1.74,-6.5)--(9.84,-6.32)--cycle, linewidth(1) + zzttqq); 
         /* draw figures */
        draw((1.08,-0.46)--(-1.74,-6.5), linewidth(1) + zzttqq); 
        draw((-1.74,-6.5)--(9.84,-6.32), linewidth(1) + zzttqq); 
        draw((9.84,-6.32)--(1.08,-0.46), linewidth(1) + zzttqq); 
        draw((-1.74,-6.5)--(2.0760281866790997,-4.028311415249227), linewidth(1)); 
        draw((2.0760281866790997,-4.028311415249227)--(3.4171770574334235,-2.0234540589680203), linewidth(1)); 
        draw((2.0760281866790997,-4.028311415249227)--(-0.10957375257483397,-3.0078813707631187), linewidth(1)); 
        draw((2.0760281866790997,-4.028311415249227)--(2.113517140377381,-6.440100769838693), linewidth(1)); 
        draw((2.0760281866790997,-4.028311415249227)--(9.84,-6.32), linewidth(1)); 
        draw((2.0760281866790997,-4.028311415249227)--(2.746488258178991,-6.430261840546442), linewidth(1)); 
        draw((2.0760281866790997,-4.028311415249227)--(1.08,-0.46), linewidth(1)); 
        draw(circle((2.4112582224290486,-5.229286627897834), 1.246884372452104), linewidth(1)); 
        draw((-0.10957375257483397,-3.0078813707631187)--(1.4523625510959888,-4.432267349734503), linewidth(1) + blue); 
        draw((1.4523625510959888,-4.432267349734503)--(2.113517140377381,-6.440100769838693), linewidth(1) + blue); 
        draw((-0.10957375257483397,-3.0078813707631187)--(2.113517140377381,-6.440100769838693), linewidth(1) + blue);
        draw((2.113517140377381,-6.440100769838693)--(3.34491881817574,-4.402849376671095), linewidth(1) + qqwuqq); 
        draw((3.34491881817574,-4.402849376671095)--(3.4171770574334235,-2.0234540589680203), linewidth(1) + qqwuqq); 
        draw((2.113517140377381,-6.440100769838693)--(3.4171770574334235,-2.0234540589680203), linewidth(1) + qqwuqq);
        draw((3.34491881817574,-4.402849376671095)--(2.746488258178991,-6.430261840546442), linewidth(1) + ccqqqq); 
        draw((1.4523625510959888,-4.432267349734503)--(3.34491881817574,-4.402849376671095), linewidth(1)); 
        draw((1.4523625510959888,-4.432267349734503)--(2.746488258178991,-6.430261840546442), linewidth(1) + ccqqqq); 
        draw(circle((-2.0422514878068454,-6.695772021481385), 4.163625905375834), linewidth(1)); 
        draw(circle((-1.152018529372594,-3.0754901031566897), 4.688744837241597), linewidth(1)); 
        draw((-5.605528239608501,-4.541975496740346)--(3.4171770574334235,-2.0234540589680203), linewidth(1) + linetype("2 2") + wwwwww); 
        draw((3.34491881817574,-4.402849376671095)--(-5.605528239608501,-4.541975496740346), linewidth(1) + linetype("2 2") + wwwwww); 
        draw((2.113517140377381,-6.440100769838693)--(-5.605528239608501,-4.541975496740346), linewidth(1) + linetype("2 2") + wwwwww); 
         /* dots and labels */
        dot((1.08,-0.46),dotstyle); 
        label("$A$", (1.16,-0.26), NE * labelscalefactor); 
        dot((-1.74,-6.5),dotstyle); 
        label("$B$", (-1.66,-6.3), SW * labelscalefactor * 5); 
        dot((9.84,-6.32),dotstyle); 
        label("$C$", (9.92,-6.12), NE * labelscalefactor); 
        dot((2.0760281866790997,-4.028311415249227),linewidth(4pt) + dotstyle); 
        label("$I$", (2.16,-3.86), N * labelscalefactor); 
        dot((2.113517140377381,-6.440100769838693),linewidth(4pt) + dotstyle); 
        label("$D$", (2.2,-6.28), S * labelscalefactor * 4); 
        dot((3.4171770574334235,-2.0234540589680203),linewidth(4pt) + dotstyle); 
        label("$E$", (3.5,-1.86), NE * labelscalefactor); 
        dot((-0.10957375257483397,-3.0078813707631187),linewidth(4pt) + dotstyle); 
        label("$F$", (-0.02,-2.84), NW * labelscalefactor); 
        dot((2.746488258178991,-6.430261840546442),linewidth(4pt) + dotstyle); 
        label("$K$", (2.82,-6.28), SE * labelscalefactor * 3); 
        dot((1.4523625510959888,-4.432267349734503),linewidth(4pt) + dotstyle); 
        label("$P$", (1.54,-4.28), NW * labelscalefactor); 
        dot((3.34491881817574,-4.402849376671095),linewidth(4pt) + dotstyle); 
        label("$Q$", (3.42,-4.24), NE * labelscalefactor); 
        dot((-5.605528239608501,-4.541975496740346),linewidth(4pt) + dotstyle); 
        label("$T$", (-5.92,-4.58), NW * labelscalefactor); 
        clip((xmin,ymin)--(xmin,ymax)--(xmax,ymax)--(xmax,ymin)--cycle); 
         /* end of picture */
    \end{asy}
\end{center}

Karena $PQ \parallel BC$ maka $\angle TDB = \angle DTP = \angle DFP$. Karena $FP=PD$ maka $\angle TDB = \angle PDF$. Ini menyebabkan $\angle PDT = \angle PDF + \angle FDT = \angle TDB + \angle FDT = \angle FDB$.

Di lain sisi, karena $FP=PD$ dan $FB=BD$ maka $BP$ melewati titik pusat lingkaran $(FPD)$. Lalu, karena $IK$ diameter $\omega$, maka $PK \perp BP$. Padahal kita juga punya $FD \perp BP$. Ini berarti $PK \parallel FD$ dan secara khusus $PK$ menyinggung lingkaran $(FPD)$ di $P$. Oleh karena itu, kita punya $\angle PKD = \angle FDB$.

Dari fakta-fakta tersebut, akan didapatkan $\angle PQD = \angle PKD = \angle FDB = \angle PDT$, yang membuktikan bahwa $TD$ menyinggung $\omega$ di $D$, seperti yang ingin dibuktikan.

    
\end{proof}